\section*{10月月考复习卷}
\begin{enumerate}
    \item 若直线 $l$ 过点 $A\left( {5, - 4}\right)$ ,且它在两坐标轴上的截距互为相反数,则直线 $l$ 的方程为\blkda[4]{$y = -\dfrac{4}{5}x$或$y = x-9$}.
    \item 若 $\bigtriangleup {ABC}$ 的三个顶点的坐标分别是 $A\left( {3, - 1}\right) ,B\left( {-5,3}\right) ,C\left( {1, - 5}\right)$ ,则 $\angle A$ 的角平分线所在的直线方程是\blkda[3]{$x-3y-6 = 0$}.
    \item 若三条直线 ${l}_{1} : x - y = 0,{l}_{2} : x + y - 2 = 0$ 和 ${l}_{3} : {5x} - {ky} - {15} = 0$ 能围成一个三角形, 则实数 $k$ 的取值范围是\blkda[4]{$\{k | k \neq \pm 5, k \neq -10\}$}.
    \begin{jieda}
        直线$l_1, l_2$交点为$(1, 1)$，当直线$l_3$过该点或与另两条直线平行时不能围成三角形，否则均可以.
    \end{jieda}
    \item 直线 ${l}_{1} : x + \left( {m + 1}\right) y - {2m} - 2 = 0$ 与直线 ${l}_{2} : \left( {m + 1}\right) x - y - {2m} - 2 = 0$ 相交于点 $P$ ,对任意实数 $m$ ,直线 ${l}_{1},{l}_{2}$ 分别恒过定点 $A,B$ ,则 $\left| {PA}\right|  + \left| {PB}\right|$ 的最大值为\blkda{$4$}.
    \begin{jieda}
        因为 ${l}_{1} : x + \left( {m + 1}\right) y - {2m} - 2 = 0$ ,即 $x + y - 2 + m\left( {y - 2}\right)  = 0$，由 $\left\{  \begin{array}{l} x + y - 2 = 0 \\  y - 2 = 0 \end{array}\right.$ ,解得 $\left\{  \begin{array}{l} x = 0 \\  y = 2 \end{array}\right.$ ,所以直线 ${l}_{1}$ 过定点 $A\left( {0,2}\right)$ ;
        同理可得直线 ${l}_{2}$ 过定点 $B\left( {2,0}\right)$ ；
        
        又因为 $1 \times  \left( {m + 1}\right)  + \left( {-1}\right)  \times  \left( {m + 1}\right)  = 0$ ,所以 ${l}_{1} \bot  {l}_{2}$ ,
        即有 ${PA}\bot {PB}$ ,所以 ${\left| PA\right| }^{2} + {\left| PB\right| }^{2} = {\left| AB\right| }^{2} = 8$
        
        所以 $\left| {PA}\right|  + \left| {PB}\right|  \leq  \sqrt{2\left( {{\left| PA\right| }^{2} + {\left| PB\right| }^{2}}\right) } = 4$ ,
        当且仅当 $\left| {PA}\right|  = \left| {PB}\right|  = 2$ 时,取等号.
        所以 $\left| {PA}\right|  + \left| {PB}\right|$ 的最大值为 4 .
    \end{jieda}
    \item 若入射光线所在的直线为 ${l}_{1} : x - y + 1 = 0$ 射入,被镜面直线 ${l}_{2} : {2x} + y - 4 = 0$ 反射,则反射光线所在的直线 $l$ 的方程是\blkda[3]{$x-7y+13=0$}.
    \begin{jieda}
        $l_1, l_2$的交点为$P(1,2)$
        
        在$l_1$上取点$M(0, 1)$，设$M$关于$l_2$的对称点为$M'(a, b)$，则$\vv{MM'} = (a, b - 1)$和$l_2$的方向向量$\vv{d} = (-1, 2)$垂直，且$MM'$的中点$N\left(\dfrac{a}{2}, \dfrac{b+1}{2}\right)$在直线$l_2$上，故$\left\{\begin{array}{l}
             -a+2(b-1) = 0  \\
             a + \dfrac{b+1}{2} - 4 = 0
        \end{array}\right.$，故$M'\left(\dfrac{12}{5}, \dfrac{11}{5}\right)$

        直线$l$即直线$M'P$，因此直线$l: x-7y+13=0$
    \end{jieda}
    \item 已知直线 ${l}_{1}$ 过点 $M$ ,将直线 ${l}_{1}$ 绕点 $M$ 沿逆时针方向旋转 $a\left( {{0}^{ \circ  } < a < {90}^{ \circ  }}\right)$ 角,得到 的直线 ${l}_{2}$ 的方程是 $x + y - 2 = 0$ ,将直线 ${l}_{2}$ 绕点 $M$ 沿顺时针方向继续旋转 $\left( {{90}^{ \circ  } - a}\right)$ 角,得到的直线 ${l}_{3}$ 的方程是 ${2x} - y - 1 = 0$ ,则直线 ${l}_{1}$ 的方程是\blkda[3]{$2x+y-3 = 0$}.
    \begin{jieda}
        $l_2$和$l_3$的交点为$M$，故$M(1, 1)$

        $k_2 = -1, k_3 = 2, tan(\ang{90} - a) = \left|\dfrac{k_3-k_2}{1+k_2k_3}\right| = 3$，故$tan(a) = \dfrac{1}{3}$
        
        $\left|\dfrac{k_1-k_2}{1+k_1k_2}\right| = \dfrac{1}{3}$，解得$k_1 = -2$，因此直线 ${l}_{1}$ 的方程是$2x+y-3 = 0$
    \end{jieda}
    \item 设 $M\left( {{x}_{1},{y}_{1}}\right) ,N\left( {{x}_{2},{y}_{2}}\right)$ 为不同的两点,直线 $l : {ax} + {by} + c = 0,\delta  = \dfrac{a{x}_{1} + b{y}_{1} + c}{a{x}_{2} + b{y}_{2} + c}$ , 以下命题中正确的序号为\dotfill{(\kh{D})}
    \begin{dingautolist}{172}
        \item 存在实数 $\delta$ ,使得点 $N$ 在直线 $l$ 上
        \item 若 $\delta  = 1$ ,则过 $M$ 、 $N$ 的直线与直线 $l$ 平行
         \item 若 $\delta  =  - 1$ ,则直线 $l$ 经过 ${MN}$ 的中点
        \item 若 $0 < \delta  < 1$ ,则点 $M$ 、 $N$ 在直线 $l$ 的同侧且直线 $l$ 与线段 ${MN}$ 的反向延长线相交
    \end{dingautolist}
    \xx{\ding{172}\ding{173}\ding{174}}{\ding{173}\ding{175}}{\ding{174}\ding{175}}{\ding{173}\ding{174}\ding{175}}

    
    \begin{jieda}
        \begin{dingautolist}{172}
            \item 点 $N$ 在直线 $l$ 上$\Longleftrightarrow $分母为0，故不可能.
            \item $\delta  = 1 \Longleftrightarrow a(x_1-x_2)+b(y_1-y_2) = 0$，故过 $M,N$ 两点的直线与直线 $l$ 平行（法向量相同），但点 $N$ 不在直线 $l$ 上，故必然是平行关系.
            \item $\delta  = -1 \Longleftrightarrow a(x_1+x_2)+b(y_1+y_2) + 2c = 0 \Longleftrightarrow a \cdot \dfrac{x_1+x_2}{2} + b \cdot \dfrac{y_1+y_2}{2} + c = 0$，故直线 $l$ 经过线段 ${MN}$ 的中点.
            \item 当 $0 < \delta  < 1$ ,得到 $0 < \dfrac{a{x}_{1} + b{y}_{1} + c}{a{x}_{2} + b{y}_{2} + c} < 1$ ,即 $\left( {a{x}_{1} + b{y}_{1} + c}\right) \left( {a{x}_{2} + b{y}_{2} + c}\right)  > 0$ ,所以得到点 $M, N$ 在直线 $l$ 的同侧,且 $\left| {a{x}_{1} + b{y}_{1} + c}\right|  < \left| {a{x}_{2} + b{y}_{2} + c}\right|$ ,得到点 $M$ 与点 $N$ 到直线 $l$ 的距离不等，所以直线 $l$ 与线段 ${MN}$ 的反向延长线相交
        \end{dingautolist}
    \end{jieda}




    \item 已知 $\bigtriangleup  {ABC}$ 的三个顶点 $A\left( {m,n}\right)$ 、 $B\left( {2,1}\right)$ 、 $C\left( {-2,3}\right)$ .
    \begin{enumerate}
        \item 求 ${BC}$ 边所在直线的方程
        \item ${BC}$ 边上中线 ${AD}$ 的方程为 ${2x} - {3y} + 6 = 0$ ,且 ${S}_{\bigtriangleup {ABC}} = 7$ ,求点 $A$ 的坐标
    \end{enumerate}
    \begin{jieda}
        \begin{enumerate}
            \item 因为 $B\left( {2,1}\right) \text{、}C\left( {-2,3}\right)$ ,所以 ${BC}$ 边所在直线的方程为：$\dfrac{y - 1}{3 - 1} = \dfrac{x - 2}{-2 - 2} \Longrightarrow  x + {2y} - 4 = 0$
            
            \item ${BC}$ 边上中线 ${AD}$ 的方程为 ${2x} - {3y} + 6 = 0$ ,所以有 ${2m} - {3n} + 6 = 0$ 

            $\vv{CB} = (4, -2), \vv{CA} = (m+2, n-3)$，$S = \dfrac{1}{2} \times |4(n-3) - (-2)(m+2)| = 7 \Longrightarrow \left| {m + {2n} - 4}\right|  = 7$
            
            因此有 $\left\{  \begin{array}{l} m + {2n} - 4 = 7 \\  {2m} - {3n} + 6 = 0 \end{array}\right.$ 或 $\left\{  \begin{array}{l} m + {2n} - 4 =  - 7 \\  {2m} - {3n} + 6 = 0 \end{array}\right.$ ,解得： $\left\{  \begin{array}{l} m = 3 \\  n = 4 \end{array}\right.$ 或 $\left\{  \begin{array}{l} m =  - 3 \\  n = 0 \end{array}\right.$，所以点 $A$ 的坐标为:$(3,4)$或$(-3,0)$.
        \end{enumerate}
    \end{jieda}

    
    \item 过点 $P\left( {2,1}\right)$ 的直线 $l$ 分别交 $y = x\left( {x \geq  0}\right)$ 与 $y =  - x\left( {x \geq  0}\right)$ 于 $A\text{ 、 }B$ 两点.
    \begin{enumerate}
        \item 若直线 $l$ 的倾斜角为 $\dfrac{\pi }{4}$ ,求直线 $l$ 的一般式方程.
        \item 当 $\left| {PA}\right|  \cdot  \left| {PB}\right|$ 最小时,求直线 $l$ 的方程
        \item 已知 $O$ 为坐标原点,设 $\triangle{AOB}$ 的面积为 ${S}$ ,讨论这样的直线 $l$ 的条数
    \end{enumerate}
    \begin{jieda}
        \begin{enumerate}
            \item 若直线 $l$ 的倾斜角为 $\dfrac{\pi }{4}$ ,则直线 $l$ 的方程为 $y - 1 = 1\left( {x - 2}\right)$ ,即 $x - y - 1 = 0$ 
            
            \item 
            \begin{enumerate}[1.]
                \item 方法1 
                \begin{dingautolist}{172}
                    \item 当直线 $l$ 的斜率不存在时, $\left| {PA}\right| \left| {PB}\right|  = 3$
                    \item 当直线 $l$ 的斜率存在时,设直线 $l : y - 1 = k\left( {x - 2}\right) ,k \in  \left( {-\infty , - 1}\right)  \cup  \left( {1, + \infty }\right)$ ,
                    $\left\{  {\begin{matrix} y = x \\  y - 1 = k\left( {x - 2}\right)  \end{matrix}\text{ 得 }A\left( {\dfrac{{2k} - 1}{k - 1},\dfrac{{2k} - 1}{k - 1}}\right) ,\left\{  {\begin{matrix} y =  - x \\  y - 1 = k\left( {x - 2}\right)  \end{matrix}\text{ 得 }B\left( {\dfrac{{2k} - 1}{k + 1},\dfrac{1 - {2k}}{k + 1}}\right) ,}\right. }\right.$
    
                    $\vv{PA} = \left( {\dfrac{1}{k - 1},\dfrac{k}{k - 1}}\right) ,\vv{PB} = \left( {\dfrac{-3}{k + 1},\dfrac{-{3k}}{k + 1}}\right)$ ,且 $\vv{PA},\vv{PB}$ 反向，所以 $\left| {PA}\right| \left| {PB}\right|  = \left| {\vv{PA} \cdot  \vv{PB}}\right|  = \dfrac{3 + 3{k}^{2}}{\left| {k}^{2} - 1\right| } = 3 + \dfrac{6}{\left| {k}^{2} - 1\right| } > 3$ 
                \end{dingautolist}
                综上所述: $\left| {PA}\right|  \cdot  \left| {PB}\right|$ 的最小值为 3,此时直线 $l$ 的斜率不存在,直线方程为 $x = 2$ .
                \item 方法2
                
                设直线$l$与$y = x$的夹角为$\theta$，$\theta \in \left(0, \dfrac{\pi}{2}\right)$则$|PA| = \dfrac{1}{\sqrt{2}} \times \dfrac{1}{sin\theta}$，$|PB| = \dfrac{3}{\sqrt{2}} \times \dfrac{1}{cos\theta}$，故$|PA|\cdot|PB| = \dfrac{3}{sin2\theta} \leq 3$，当$\theta = \dfrac{\pi}{4}$时取得最小值$3$，此时直线方程为$x = 2$.
            \end{enumerate}
            
            \item 
            \begin{dingautolist}{172}
                \item 当直线斜率不存在时, $A\left( {2,2}\right) ,B\left( {2, - 2}\right) ,S = \dfrac{1}{2}\left| {OA}\right|  \cdot  \left| {OB}\right|  = 4$
                \item 当直线斜率存在时, $k \in  \left( {-\infty , - 1}\right)  \cup  \left( {1, + \infty }\right) ,A\left( {\dfrac{{2k} - 1}{k - 1},\dfrac{{2k} - 1}{k - 1}}\right) ,B\left( {\dfrac{{2k} - 1}{k + 1},\dfrac{1 - {2k}}{k + 1}}\right)$ ,
                
                $S = \dfrac{1}{2} \left|\dfrac{{2k} - 1}{k - 1} \cdot \dfrac{1 - {2k}}{k + 1} - \dfrac{{2k} - 1}{k - 1} \cdot \dfrac{{2k} - 1}{k + 1}\right| = \dfrac{{\left( 2k - 1\right) }^{2}}{{k}^{2} - 1}$，即 $\left( {4 - S}\right) {k}^{2} - {4k} + S + 1 = 0$
                \begin{enumerate}[1.]
                    \item 方法1 
                    
                    当 $S = 4$ 时,方程有 1 解,此时 $k = \dfrac{5}{4}$
                    
                    当 $S \neq  4$ 时, $\Delta  = {16} - 4\left( {4 - S}\right) \left( {S + 1}\right)  = {4S}\left( {S - 3}\right)$ 
                    
                    当 $S < 3$ 时, $\Delta  < 0$ ,方程无解
                    
                    当 $S = 3$ 时, $\Delta  = 0,k = 2$ ,方程有 1 解
                    
                    当 $4 > S > 3$ 时, $\Delta  > 0,f\left( k\right)  = \left( {4 - S}\right) {k}^{2} - {4k} + S + 1$，开口向上，对称轴 $\dfrac{2}{4 - S} > 2$ ,且 $f\left( 1\right)  = 1 > 0$ ,方程有两个大于 1 的解.
                    
                    当 $S > 4$ 时, $\Delta  > 0,f\left( k\right)  = \left( {4 - S}\right) {k}^{2} - {4k} + S + 1$ 开口向下, $f\left( 1\right)  = 1 > 0,f\left( {-1}\right)  = 9 > 0$ ,方程有 1 个大于 1 的解, 一个小于 -1 的解.
                    \item 方法2

                    $S = \dfrac{(2k-1)^2}{k^2-1}$
                    
                    令$t = 2k-1$，$t \in (-\infty, -3) \cup (1, +\infty)$，$t$与$k$一一对应，则$S = \dfrac{t^2}{\left(\dfrac{t+1}{2}\right)^2-1} = \dfrac{4t^2}{t^2+2t-3} = \dfrac{4}{1+\dfrac{2}{t}-\dfrac{3}{t^2}}$，即$1+\dfrac{2}{t}-\dfrac{3}{t^2} = \dfrac{4}{S}$

                    令$m = \dfrac{1}{t}$，$m \in \left(-\dfrac{1}{3}, 0\right) \cup (0, 1)$，$t$与$m$一一对应，则$-3m^2+2m+1 = \dfrac{4}{S}$
                    
                    $g(m) = -3m^2+2m+1$在$\left(-\infty, \dfrac{1}{3}\right)$严格增，在$\left(\dfrac{1}{3}, +\infty\right)$严格减
                    
                    数形结合可知当$\dfrac{4}{S} \in [0, 1)\cup \left(1, \dfrac{4}{3}\right)$时原方程有两解，当$\dfrac{4}{S} \in \left\{1, \dfrac{4}{3}\right\}$时，原方程只有一解，即$S \in (3, 4) \cup (4, +\infty)$时原方程有两解，当$S \in \{3,4\}$原方程只有一解.
                \end{enumerate}
            \end{dingautolist}

            综上所述：当 $S < 3$ 时, 0 条；当 $S = 3$ 时, 1 条；当 $S > 3$ 时, 2 条.
        \end{enumerate}
    \end{jieda}
    \item 曲线 ${x}^{2} + {y}^{2} = 4\left| x\right|  + 4\left| y\right|$ 的周长为\blkda[3]{$8\sqrt{2}\pi$}
    \begin{jieda}
        在第一象限及$x$轴,$y$轴正方向上，方程可化为${x}^{2} + {y}^{2} = 4x+4y$，即$(x-2)^2+(y-2)^2 = 8$，计算可知该段周长为$2\sqrt{2}\pi$

        若$(x_0, y_0)$在曲线上，则$(x_0, -y_0), (-x_0, y_0), (-x_0, -y_0)$均在曲线上，因此曲线关于$x$轴，$y$轴对称，且关于原点中心对称，故周长为$8\sqrt{2}\pi$
    \end{jieda}
    \item 已知点 $A\left( {0,1}\right)$ 、 $B\left( {4,t}\right)$ . 是否存在实数 $t$ ,满足过 $A$ 、 $B$ 两点作与 $x$ 轴相切的圆有且只有一个? 若存在满足条件的圆, 求出这个圆的方程; 若不存在满足条件的圆, 说明理由.
    \begin{jieda}
        设圆心为$(a, b)$，根据题意可知半径为$b$，故圆的标准方程为$(x-a)^2+(y-b)^2 = b^2$
        
        代入点$A, B$的坐标可知$\left\{\begin{array}{l}
             a^2+(1-b)^2 = b^2 \\
             (4-a)^2+(t-b)^2 = b^2
        \end{array}\right.$

        上式化简可得$a^2+1=2b$，两式做差可得$4(4-2a)+(t+1-2b)(t-1) = 0$，消元得$4(4-2a)+(t-a^2)(t-1) = 0$，整理成$a$为主元的形式：$(1-t)a^2 - 8a + 16+t^2-t = 0$

        存在实数 $t$ ,满足过 $A$ 、 $B$ 两点作与 $x$ 轴相切的圆有且只有一个，即存在实数 $t$使得关于$a$的方程$(1-t)a^2 - 8a + 16+t^2-t = 0$只有一个实根，即$t = 1$或$\Delta = 64 - 4(1-t)(16+t^2-t) = 0 (t \neq 1)$

        $t \neq 1$时，$64 - 4(1-t)(16+t^2-t) = 0$即$\dfrac{16}{1-t} = t^2-t+16$，结合图像可知$y = \dfrac{16}{1-t}$与$y = t^2-t+16$只有一个交点$(0, 16)$.

        故$t \in \{0, 1\}$时满足要求.

        \begin{dingautolist}{172}
            \item $t=1,(x-2)^2+\left(y-\dfrac{5}{2}\right)^2=\dfrac{25}{4}$
            \item $t=0,(x-4)^2+\left(y-\dfrac{17}{2}\right)^2=\dfrac{289}{4}$
        \end{dingautolist}
    \end{jieda}

    \item 过点 $A\left( {1, - 3}\right)$ 的圆 ${x}^{2} + {y}^{2} = {10}$ 的切线方程是 \blkda[4]{$x-3y-10=0$}.
    \begin{jieda}
        $A$在圆上，为切点，故切线方程为$x\times1+y\times3 = 10$，即$x-3y-10=0$.
    \end{jieda}

    \item 圆 ${x}^{2} + {y}^{2} = 9$ 的过点 $M\left( {3,4}\right)$ 的切线方程为\blkda[6]{$7x-24y+75=0$或$x=3$}.
    \begin{jieda}
        $M$在圆外
        \begin{dingautolist}{172}
            \item 当斜率不存在时，过$M$的直线为$x = 3$，满足要求
            \item 当斜率存在时，设切线为$y = k(x-3)+4$，故原点到切线的距离$d = \dfrac{|4-3k|}{\sqrt{1+k^2}} = 3$，解得$k = \dfrac{7}{24}$，故切线方程为$7x-24y+75=0$
        \end{dingautolist}
    \end{jieda}
    \item 与圆 ${x}^{2} + {y}^{2} + {8x} + {6y} = 0$ 相切且在两坐标轴上的截距相等的直线方程为\blkda[11]{$x+y+7+5\sqrt{2}=0$或$x+y+7-5\sqrt{2}=0$或$4x+3y=0$}.
    \begin{jieda}
        ${x}^{2} + {y}^{2} + {8x} + {6y} = 0 \Longleftrightarrow (x+4)^2+(y+3)^2 = 25$.
        
        \begin{dingautolist}{172}
            \item 截距为$0$时，设直线方程为$y = kx$，则$d = \dfrac{|3k-4|}{\sqrt{1+k^2}} = 5$，解得$k = -\dfrac{4}{3}$，故直线为$4x+3y = 0$
            \item 截距不为$0$时，设直线方程为$\dfrac{x}{a} + \dfrac{y}{a} = 1$，即$x+y = a (a \neq 0)$，则$d = \dfrac{|-3-4-a|}{\sqrt{2}} = 5$，解得$a = -7 \pm 5\sqrt{2}$，故直线方程为$x+y+7 \pm 5\sqrt{2} = 0$
        \end{dingautolist}
    \end{jieda}

    \item 已知直线 ${l}_{1} : x - {8y} + {30} = 0$ 与直线 ${l}_{2} : x + {5y} - {22} = 0$ 的交点为点 $A$ ,直线 ${l}_{3}$ 过点 $A$ ,且直线 ${l}_{3}$ 被圆 ${x}^{2} + {y}^{2} - {2x} + {2y} - {14} = 0$ 所截得的线段的长为 $2\sqrt{3}$ ,直线 ${l}_{3}$ 的方程为\blkda[8]{$2 x-3 y+8=0$ 和 $3 x+2 y-14=0$ }.
    \begin{jieda}
        由 $\left\{\begin{array}{l}x-8 y+30=0, \\ x+5 y-22=0,\end{array}\right.$ 得 $\left\{\begin{array}{l}x=2, \\ y=4 .\end{array}\right.$ 
        
        易知斜率不存在时不满足要求，故设直线 $l_3$ 的方程为 $y-4=k(x-2)$．因为圆 $(x-1)^2+(y+1)^2 =16$ 的圆心为 $(1,-1)$ ，半径 $r=4$ ，所以圆心到直线 $l$ 的距离 $d=\dfrac{|5-k|}{\sqrt{1+k^2}}$ ．由题设，得 $\dfrac{k^2-10 k+25}{1+k^2}+ (\sqrt{3})^2=16$ ．解得 $k_1=\dfrac{2}{3}, k_2=-\dfrac{3}{2}$ ．因此，直线 $l_3$ 的方程为 $2 x-3 y+8=0$ 和 $3 x+2 y-14=0$ ．
    \end{jieda}
    \item 若关于 $x$ 的方程 ${kx} = \sqrt{-{x}^{2} + {4x} - 3} + 1$ 有两个不同的实数解,则实数 $k$ 的取值范围是\blkda[3]{$\left[1,\dfrac{4}{3}\right)$}.
    \begin{jieda}
        方程 $k x=\sqrt{-x^2+4 x-3}+1 \Longleftrightarrow k x-1=\sqrt{1-(x-2)^2}$ ，令 $y=k x-1, y=\sqrt{1-(x-2)^2}$ ，由 $y=\sqrt{1-(x-2)^2}$ ，得 $(x-2)^2+y^2=1$ ，因此 $y=\sqrt{1-(x-2)^2}$ 表示以点 $C(2,0)$ 为圆心， 1 为半径的上半圆（含端点）， $y=k x-1$ 表示纵截距为 -1 ，斜率为 $k$ 的直线 $l$ ，依题意，直线 $l$ 与半圆 $C$ 有两个交点，
        在同一坐标系内作出直线 $l$ 与半圆 $C$ 
        \begin{center}
            \includegraphics[width=0.25\linewidth]{bodyxb//src/zhixianheyuan1.png}
        \end{center}
        
        如图,当直线 $l$ 与半圆 $C$ 相切时，$\dfrac{|2 k-1|}{\sqrt{k^2+1}}=1$ ，显然 $k>0$ ，解得 $k=\dfrac{4}{3}$ ，当直线 $l$ 过点 $(1,0)$ 时，$k=1$ ，此时直线 $l$ 与半圆 $C$ 有两个交点，当 $1 \leq k<\dfrac{4}{3}$ 时，直线 $l$ 与半圆 $C$ 有两个交点，即方程 $k x=\sqrt{-x^2+4 x-3}+1$ 有两个不同的实数解，所以实数 $k$ 的取值范围是 $1 \leq k<\dfrac{4}{3}$ ．
    \end{jieda}
    \item 圆 ${x}^{2} + {y}^{2} + {2x} + {4y} - 3 = 0$ 上到直线 $x + y + 1 = 0$ 的距离等于 $\sqrt{2}$ 的点有 \blkda[2]{3}个
    \item 已知圆 $M : {x}^{2} + {y}^{2} - {2ay} + {a}^{2} - 4 = 0$ 与圆 $N : {x}^{2} + {y}^{2} - {4x} + 3 = 0$ 有 4 条公切线, 则 $a$ 的取值范围为\blkda[5]{$(-\infty,-\sqrt{5})\cup（\sqrt{5},+\infty)$}.
    \begin{jieda}
        由题可得，圆 $M: x^2+(y-a)^2=4$ ，圆心为 $(0, a)$ ，半径为 2 ，圆 $N:(x-2)^2+y^2=1$ ，圆心为 $(2,0)$ ，半径为 1 ，因为两圆有 4 条公切线，所以两圆外离，故圆心距 $|M N|=\sqrt{2^2+a^2}>3$ ，解得 $a \in(-\infty,-\sqrt{5}) \cup(\sqrt{5},+\infty)$.
    \end{jieda}
    \item 圆系: ${C}_{k} : {x}^{2} + {y}^{2} + {2kx} + \left( {{4k} + {10}}\right) y + {10k} + {20} = 0\left( {k \neq   - 1}\right)$ 中任两圆之间的关系是
    \dotfill{(\kh{D})}
    \xx{相交}{相离}{外切}{内切或外切}
    \item 已知直线 $x + {2y} - 3 = 0$ 和圆 ${x}^{2} + {y}^{2} + x - {2ay} + a = 0$ 相交于 $A$ 、 $B$ 两点,且以 ${AB}$ 为直径的圆经过坐标原点,求实数 $a$ 的值.
    \begin{jieda}
        过$A,B$两点的圆系方程为$x^2+y^2+(1-\lambda)x-2(a+\lambda)y+(a+3\lambda) = 0$
        \begin{enumerate}
            \item 方法1
            
                半径$r = \dfrac{1}{2}\sqrt{(1-\lambda)^2+4(a+\lambda)^2-4(a+3\lambda)} = \dfrac{1}{2}\sqrt{5\lambda^2+(8a-14)\lambda + (4a^2-4a+1)}$，当$\lambda = \dfrac{7-4a}{5}$时，半径取得最小值，即以$A, B$为直径的圆，又因为该圆过原点，所以有$a + 3\lambda = 0$，两式联立$\left\{\begin{array}{l}
                 \lambda = \dfrac{7-4a}{5}  \\
                 a + 3\lambda = 0
            \end{array}\right.$，解得$\left\{\begin{array}{l}
                 \lambda = -1  \\
                 a = 3
            \end{array}\right.$.
            \item 方法2
            
            $x^2+y^2+(1-\lambda)x-2(a+\lambda)y+(a+3\lambda) = 0$应满足两个约束：过原点，圆心在$x + {2y} - 3 = 0$上

            满足过原点时方程变为$x^2+y^2+(1-\lambda)x+4\lambda y = 0$，圆心为$\left(\dfrac{\lambda-1}{2}, -2\lambda\right)$，代入直线方程可得$\lambda = -1$，故$a=3$.
        \end{enumerate}
        
    \end{jieda}

    \item 在平面直角坐标系 ${xOy}$ 中,已知两点 $S\left( {4,0}\right) ,T\left( {1,0}\right)$ ,动点 $P$ 满足 $\left| {PS}\right|  = 2\left| {PT}\right|$ , 设点 $P$ 的轨迹为 $C$ . 如图,动直线 $l$ 与曲线 $C$ 交于不同的两点 $A,B$ ( $A,B$ 均在 $x$ 轴上方),且 $\angle {ATO} + \angle {BTO} = {180}^{ \circ  }$ .
    \begin{center}
        \includegraphics[width=0.25\linewidth]{bodyxb//src/bx01ch2.2-20.png}
    \end{center}
    
    \begin{enumerate}
        \item 求曲线 $C$ 的方程；
        \item 当 ${A}$ 为曲线 $C$ 与 $y$ 轴正半轴的交点时,求直线 $l$ 的方程;
        \item 是否存在一个定点,使得直线 $l$ 始终经过此定点？若存在,求出定点的坐标；若不存在,请说明理由. 
    \end{enumerate}
    \begin{jieda}
        \begin{enumerate}
            \item 设 ${P}({x}, {y})$ ，由 $|P S|=2|P T|$ 得 $\sqrt{(x-4)^2+y^2}=2 \sqrt{(x-1)^2+y^2}$ ，化简得 $x^2+y^2=4$ ，则曲线 $C$ 的方程为 $x^2+y^2=4$ ；
            \item 由题意知 $A(0,2)$ ，设 ${B}\left({x}_2, {y}_2\right)$ ，
            依题意可知直线 $l$ 的斜率存在，设直线 $l$ 的方程为 $y=k x+2$ ，
            由 $\angle A T O+\angle B T O=180^{\circ}$ ，得 $k_{A T}+k_{B T}=0$ ，
            则 $\left\{\begin{array}{l}-2+\dfrac{y_2}{x_2-1}=0 \\ x_2^2+y_2^2=4\end{array}\right.$ ，所以 $\left\{\begin{array}{l}x_2=0 \\ y_2=-2\end{array}\right.$（舍去）或 $\left\{\begin{array}{l}x_2=\dfrac{8}{5} \\ y_2=\dfrac{6}{5}\end{array}\right.$ ，即 $B\left(\dfrac{8}{5}, \dfrac{6}{5}\right)$ ，
            则 $k=\dfrac{\dfrac{6}{5}-2}{\dfrac{8}{5}-0}=-\dfrac{1}{2}$ ，
            则直线 $l$ 方程为 $y=-\dfrac{1}{2} x+2$ ；
            
            
            \item 设直线 $l$ 方程为 $y=k x+b$ ，设 ${A}\left({x}_1, {y}_1\right), {B}\left({x}_2, {y}_2\right)$ ，
            
            $\because \angle A T O+\angle B T O=180^{\circ}, \therefore k_{A T}+k_{B T}=0$
            
            $\therefore \dfrac{y_1}{x_1-1}+\dfrac{y_2}{x_2-1}=0 \Longrightarrow y_1\left(x_2-1\right)+y_2\left(x_1-1\right)=0$
             
            $\therefore \left(k x_1+b\right)\left(x_2-1\right)+\left(k x_2+b\right)\left(x_1-1\right)=0$
            
            $\therefore 2 k x_1 x_2+(b-k)\left(x_1+x_2\right)-2 b=0$
            
            
            联立方程 $\left\{\begin{array}{l}x^2+y^2=4 \\ y=k x+b\end{array}\right.$ ，得 $\left(k^2+1\right) x^2+2 k b x+b^2-4=0$ ，因此$\left\{\begin{array}{l}
                 x_1+x_2=\dfrac{-2 k b}{k^2+1} \\
                 x_1 x_2=\dfrac{b^2-4}{k^2+1}
            \end{array}\right.$

            $\therefore 2 k \dfrac{b^2-4}{k^2+1}+(b-k)\left(\dfrac{-2 k b}{k^2+1}\right)-2 b=0$，即$2 k\left(b^2-4\right)-2 k b(b-k)-2 b\left(k^2+1\right)=0$

            故$b = -4k$，则直线 $l$ 始终经过此定点 $(4,0)$ ．
            
        \end{enumerate}
    \end{jieda}
    
    \item 已知圆 ${x}^{2} + {y}^{2} - {6x} - {55} = 0$ ,动圆 $M$ 经过定点 $A\left( {-3,0}\right)$ ,且与已知圆内切,求圆心 $M$ 的轨迹方程.
    \begin{jieda}
        已知圆 ${\left( x - 3\right) }^{2} + {y}^{2} = {64}$ 的圆心为 $C\left( {3,0}\right)$ ,半径 $r = 8$ . 如右图,设圆 $M$与圆 $C$ 内切于 $B$ ,则 $C$ 、 $M$ 、 $B$ 在一直线上,即 $\left| {MB}\right|  + \left| {MC}\right|  = \left| {CB}\right|  = 8$ . 又因为$\left| {MB}\right|  = \left| {MA}\right|$ ,所以 $\left| {MA}\right|  + \left| {MC}\right|  = 8$ . 因为 $8 > 6 = \left| {AC}\right|$ ,所以由椭圆定
        义,知点 $M$ 的轨迹是以 $A\left( {-3,0}\right) \text{、}C\left( {3,0}\right)$ 为焦点,长轴长为 8 的椭圆,且 $a = 4$ ,$c = 3,{b}^{2} = {a}^{2} - {c}^{2} = 7$ . 故点 $M$ 的轨迹方程是 $\dfrac{{x}^{2}}{16} + \dfrac{{y}^{2}}{7} = 1$ .
    \end{jieda}
    
    \item 求与两圆 ${\left( x + 2\right) }^{2} + {y}^{2} = 1,{x}^{2} - {4x} + {y}^{2} - {45} = 0$ 都相切的动圆圆心的轨迹方程.
    \begin{jieda}
        分两种情况，与圆1外切，与圆2内切，此时轨迹方程为
        $\dfrac{{x}^{2}}{16} + \dfrac{{y}^{2}}{12} = 1$；与两圆均内切，此时轨迹方程为$\dfrac{{x}^{2}}{9} + \dfrac{{y}^{2}}{5} = 1$
    \end{jieda}

    \item 已知椭圆的中心在坐标原点 ${O}$ ,焦点在 $x$ 轴上,直线 $y = x + 1$ 与椭圆相交于点 $P$ 和点 $Q$ ,且 ${OP} \bot  {OQ},\left| {PQ}\right|  = \dfrac{\sqrt{10}}{2}$ ,求椭圆方程.

    \begin{jieda}
        设椭圆的标准方程是$\dfrac{x^2}{a^2}+\dfrac{y^2}{b^2} = 1(a > b > 0)$，联立直线和椭圆的方程，得到方程组$\left\{\begin{array}{l}
             y = x + 1 \\
             \dfrac{x^2}{a^2}+\dfrac{y^2}{b^2} = 1
        \end{array}\right.$，消去$y$，得到$b^2x^2+a^2(x+1)^2 = a^2b^2$，整理得$(a^2+b^2)x^2+2a^2x+a^2(1-b^2) = 0$，设$P(x_1, y_2), Q(x_2, y_2)$，根据韦达定理有$\left\{\begin{array}{l}
             x_1+x_2 = -\dfrac{2a^2}{a^2+b^2} \\
             x_1x_2 = \dfrac{a^2(1-b^2)}{a^2+b^2}
        \end{array}\right.$\\
        根据${OP} \bot  {OQ}$可知$x_1x_2+y_1y_2 = 0$，即$2x_1x_2+(x_1+x_2)+1 = 0$，根据$\left| {PQ}\right|  = \dfrac{\sqrt{10}}{2}$可知$|x_1-x_2| = \dfrac{\sqrt{5}}{2}$，即$(x_1+x_2)^2-4x_1x_2 = \dfrac{5}{4}$，联立两式可得$\left\{\begin{array}{l}
             x_1+x_2 = -\dfrac{3}{2} \\
             x_1x_2 = \dfrac{1}{4}
        \end{array}\right.$，因此有$\left\{\begin{array}{l}
             -\dfrac{3}{2} = -\dfrac{2a^2}{a^2+b^2} \\
             \dfrac{1}{4} = \dfrac{a^2(1-b^2)}{a^2+b^2}
        \end{array}\right.$，解得$a^2 = 2, b^2 = \dfrac{2}{3}$，因此$\dfrac{x^2}{2}+\dfrac{3y^2}{2} = 1$
    \end{jieda}
    
    \item 以椭圆 $\dfrac{{x}^{2}}{{a}^{2}} + {y}^{2} = 1\left( {a > 1}\right)$ 上的一个端点 $B\left( {0,1}\right)$ 为直角顶点,作椭圆内接等腰直角三角形 ${ABC}$ ,这样的三角形存在吗? 若存在,最多能作几个?
    \begin{jieda}
        由题设,知 ${AB}\text{、}{BC}$ 不可能垂直或平行于 $x$ 轴,故不妨设 ${BC}$ 所在直线方程为 $y = {kx} + 1\left( {k < 0}\right)$ ,代入椭圆方程,整理得 $\left( {1 + {a}^{2}{k}^{2}}\right) {x}^{2} + {2k}{a}^{2}x = 0$ ,解得${x}_{1} = 0,{x}_{2} =  - \dfrac{2{a}^{2}k}{1 + {a}^{2}{k}^{2}}$ . 因此 $\left| {BC}\right|  = \sqrt{1 + {k}^{2}}\left| {{x}_{1} - {x}_{2}}\right|  = \dfrac{2{a}^{2}\left| k\right| \sqrt{1 + {k}^{2}}}{1 + {a}^{2}{k}^{2}}$ . 
        
        同理,用 $- \dfrac{1}{k}$ 代上式中的 $k$ ,得 $\left| {AB}\right|  = \dfrac{2{a}^{2}\sqrt{1 + {k}^{2}}}{{a}^{2} + {k}^{2}}$ .
        
        由 $\left| {BC}\right|  = \left| {AB}\right|$ ,得 $\left| k\right| \left( {{a}^{2} + {k}^{2}}\right)  = 1 + {a}^{2}{k}^{2}$ . 又 $k < 0$ ,得 $\left( {k + 1}\right) \left\lbrack  {{k}^{2} + \left( {{a}^{2} - 1}\right) k}+1\right\rbrack  = 0$ . 
        
        当 ${\left( {a}^{2} - 1\right) }^{2} - 4 < 0$ ,即 $1 < a < \sqrt{3}$ 时,方程$\left( {k + 1}\right) \left\lbrack  {{k}^{2} + \left( {{a}^{2} - 1}\right) k}+1\right\rbrack  = 0$有唯一解 $k =  - 1$ ; 
        
        当 $a = \sqrt{3}$ 时,方程 ${k}^{2} + \left( {{a}^{2} - 1}\right) k + 1 = 0$ 的解为 $k =  - 1$ ,此时方程$\left( {k + 1}\right) \left\lbrack  {{k}^{2} + \left( {{a}^{2} - 1}\right) k}+1\right\rbrack  = 0$有唯一解 $k =  - 1$ ; 
        
        当 $a > \sqrt{3}$ 时,方程$\left( {k + 1}\right) \left\lbrack  {{k}^{2} + \left( {{a}^{2} - 1}\right) k}+1\right\rbrack  = 0$有三个不同解. 
        
        综上所述,当 $1 < a \leq  \sqrt{3}$ 时,最多作出一个满足条件的三角形; 当 $a > \sqrt{3}$ 时,最多可作出三个满足条件的三角形.
    \end{jieda}
    \item 已知椭圆 $\dfrac{{x}^{2}}{9} + \dfrac{{y}^{2}}{4} = 1$ 的焦点为 ${F}_{1}$ 、 ${F}_{2}$ ,点 $P$ 为椭圆上的动点,当 $\angle {F}_{1}P{F}_{2}$ 是钝角时,则 $P$ 的横坐标的范围是\blkda[4]{$\left(-\dfrac{3\sqrt{5}}{5},\dfrac{3\sqrt{5}}{5}\right)$}.
    \item 已知椭圆 $C$ 的两个焦点分别为 ${F}_{1}\left( {-1,0}\right)$ 、 ${F}_{2}\left( {1,0}\right)$ ,短轴的两个端点分别为 ${B}_{1}\text{、}{B}_{2}$ . 
    \begin{enumerate}
        \item 若 $\bigtriangleup {F}_{1}{B}_{1}{B}_{2}$ 为等边三角形,求椭圆 $C$ 的标准方程;
        \item 若椭圆 $C$ 的短轴长为 2,过点 ${F}_{2}$ 的直线 $l$ 与椭圆 $C$ 相交于 $P$ 、 $Q$ 两点,且 $\vv{{F}_{1}P}\bot \vv{{F}_{1}Q}$ ,求直线 $l$ 的方程.
    \end{enumerate}

    \begin{jieda}
    \begin{enumerate}
        \item 由题意知, $ \bigtriangleup {F}_{1}{B}_{1}{B}_{2} $ 为等边三角形,
            所以 $ c = \sqrt{3}b $ ,又 $ c = 1 $ ,
            所以 $ b = \dfrac{\sqrt{3}}{3} $ ,
            又由 $ {a}^{2} = {b}^{2} + {c}^{2} $ ,可得 $ {a}^{2} = \dfrac{4}{3} $ ,
            故椭圆 $ C $ 的方程为 $ \dfrac{3{x}^{2}}{4} + 3{y}^{2} = 1 $ .
        \item 易知椭圆 $ C $ 的方程为 $ \dfrac{{x}^{2}}{2} + {y}^{2} = 1 $ ,
            当直线 $ l $ 的斜率不存在时,其方程为 $ x = 1 $ ,不符合题意;
            当直线 $ l $ 的斜率存在时,设直线 $ l $ 的方程为 $ y = k\left( {x - 1}\right) $ ,
            由 $ \left\{  \begin{array}{l} y = k\left( {x - 1}\right) , \\  \dfrac{{x}^{2}}{2} + {y}^{2} = 1, \end{array}\right. $
            得 $ \left( {2{k}^{2} + 1}\right) {x}^{2} - 4{k}^{2}x + 2\left( {{k}^{2} - 1}\right)  = 0 $ ,
            设 $ P\left( {{x}_{1},{y}_{1}}\right) ,Q\left( {{x}_{2},{y}_{2}}\right) $ ,
            则 $ {x}_{1} + {x}_{2} = \dfrac{4{k}^{2}}{2{k}^{2} + 1},{x}_{1}{x}_{2} = \dfrac{2\left( {{k}^{2} - 1}\right) }{2{k}^{2} + 1} $ ,
            $ \vv{{F}_{1}P} = \left( {{x}_{1} + 1,{y}_{1}}\right) , $
            $ \vv{{F}_{1}Q} = \left( {{x}_{2} + 1,{y}_{2}}\right) , $
            因为 $ \vv{{F}_{1}P} \bot  \vv{{F}_{1}Q} $ ,
            所以 $ \vv{{F}_{1}P} \cdot  \vv{{F}_{1}Q} = 0 $ ,
            即 $ \left( {{x}_{1} + 1}\right) \left( {{x}_{2} + 1}\right)  + {y}_{1}{y}_{2} = {x}_{1}{x}_{2} + \left( {{x}_{1} + {x}_{2}}\right)  + 1 + {k}^{2}\left( {{x}_{1} - }\right. $
            $ 1)\left( {{x}_{2} - 1}\right)  = \left( {{k}^{2} + 1}\right) {x}_{1}{x}_{2} - \left( {{k}^{2} - 1}\right) \left( {{x}_{1} + {x}_{2}}\right)  + {k}^{2} + 1 = $
            $ \dfrac{7{k}^{2} - 1}{2{k}^{2} + 1} = 0, $
            解得 $ {k}^{2} = \dfrac{1}{7} $ ,即 $ k =  \pm  \dfrac{\sqrt{7}}{7} $ ,
            故直线 $ l $ 的方程为 $ x + \sqrt{7}y - 1 = 0 $ 或 $ x - \sqrt{7}y - 1 = 0 $ .
    \end{enumerate}
    \end{jieda}
    \item 已知 $P\left( {{x}_{0},{y}_{0}}\right)$ 为椭圆 $\dfrac{{x}^{2}}{{a}^{2}} + \dfrac{{y}^{2}}{{b}^{2}} = 1\left( {a > b > 0}\right)$ 内且不在 $x$ 轴上的一点,直线 $l$ 过点 $P$ , 且直线 $l$ 与这个椭圆交于 $A$ 、 $B$ 两点, $A$ 、 $B$ 两点的中点恰好为点 $P$ . 求证: 直线 $l$ 的斜率 $k$ 为 $- \dfrac{{b}^{2}{x}_{0}}{{a}^{2}{y}_{0}}$ .
    \begin{jieda}
        设 $ A\text{、}B $ 两点的坐标分别为 $ \left( {{x}_{1},{y}_{1}}\right) \text{、}\left( {{x}_{2},{y}_{2}}\right) $ ,则 $\left\{\begin{array}{l}
             {b}^{2}{x}_{1}^{2} + {a}^{2}{y}_{1}^{2} = {a}^{2}{b}^{2}  \\
             {b}^{2}{x}_{2}^{2} + {a}^{2}{y}_{2}^{2} = {a}^{2}{b}^{2} 
        \end{array} \right.$ , 两式相减,得 $ {b}^{2}\left( {{x}_{1} - {x}_{2}}\right) \left( {{x}_{1} + {x}_{2}}\right)  + {a}^{2}\left( {{y}_{1} - {y}_{2}}\right) \left( {{y}_{1} + {y}_{2}}\right)  = 0 $ . 
        
        又点 $ P $ 不在 $ x $ 轴上,故 $ {x}_{1} \neq  {x}_{2} $ . 而$\left\{\begin{array}{l}
             {x}_{1} + {x}_{2} = 2{x}_{0}  \\
             {y}_{1} + {y}_{2} = 2{y}_{0}
        \end{array} \right.$
        
        所以 $ k = \dfrac{{y}_{1} - {y}_{2}}{{x}_{1} - {x}_{2}} =  - \dfrac{{b}^{2}\left( {{x}_{1} + {x}_{2}}\right) }{{a}^{2}\left( {{y}_{1} + {y}_{2}}\right) } =  - \dfrac{{b}^{2}{x}_{0}}{{a}^{2}{y}_{0}} $ .
    \end{jieda}
    
     \item 已知椭圆 $\dfrac{{x}^{2}}{4} + \dfrac{{y}^{2}}{3} = 1$ 上有不同的两点关于直线 $y = {4x} + m$ 对称,求实数 $m$ 的取值范围.
     \begin{jieda}
         由点差法，中点轨迹方程为$3x-y=0$，又因为中点在直线$y = {4x} + m$ 上，因此联立可得中点坐标为$（-m,-3m）$。因为弦中点在椭圆内，所以$\dfrac{m^2}{4}+\dfrac{9m^2}{3}<1$，从而可得$m\in \left(-\dfrac{2\sqrt{13}}{13},\dfrac{2\sqrt{13}}{13}\right)$
     \end{jieda}
    \item 若点 $P$ 是双曲线 $C : \dfrac{{x}^{2}}{4} - \dfrac{{y}^{2}}{12} = 1$ 上一点, ${F}_{1},{F}_{2}$ 分别为 $C$ 的左、右焦点, $\left| {P{F}_{1}}\right|  = 9$ , 则 $\left| {P{F}_{2}}\right|  =$ \blkda{5 或 13}.
    \begin{jieda}
        设双曲线 $C$ 的半焦距为 $c > 0$ ,则 $a = 2,b = 2\sqrt{3},c = \sqrt{{a}^{2} + {b}^{2}} = 4$
        
        因为 $\left|{P{F}_{1}}\right| - \left|{P{F}_{2}}\right| = \left| {9 - \left|{P{F}_{2}}\right|}\right|  = 4$，解得 $\left| {P{F}_{2}}\right|  = 5$ 或 $\left| {P{F}_{2}}\right|  = {13}$，经检验均符合题意.
    \end{jieda}
    \item 已知点 ${F}_{1}\text{ 、 }{F}_{2}$ 是双曲线 $\dfrac{{x}^{2}}{9} - \dfrac{{y}^{2}}{16} = 1$ 的两个焦点,点 $P$ 在此双曲线上,且满足 $\left| {P{F}_{1}}\right|  \cdot  \left| {P{F}_{2}}\right|  = {32}$ ,则 $\angle {F}_{1}P{F}_{2} =$ \blkda{$\dfrac{\pi}{2}$}.
    \begin{jieda}
        在$\triangle F_1PF_2$中列余弦定理有$|F_1F_2|^2 = |PF_1|^2+|PF_2|^2 - 2|PF_1|\cdot |PF_2| cos\angle {F}_{1}P{F}_{2} = \left(|PF_1|-|PF_2|\right)^2-2|PF_1|\cdot |PF_2| (cos\angle {F}_{1}P{F}_{2} - 1)$

        代入具体数值有$100 = 36 - 64 \times (cos\angle {F}_{1}P{F}_{2} - 1)$，故$cos\angle {F}_{1}P{F}_{2} = 0$，即$\angle {F}_{1}P{F}_{2} = \dfrac{\pi}{2}$
    \end{jieda}
    \item 已知双曲线 ${C}$ 的方程为 ${x}^{2} - \dfrac{{y}^{2}}{2} = 1$ ,试判断过点 $P\left( {1,1}\right)$ 能否作一条直线 $l$ 与双曲线交于 $A\text{ 、 }B$ 两点,且 $P$ 为 ${AB}$ 的中点,并说明理由.
    \begin{jieda}
        设双曲线上两点$A(x_1, y_1), B(x_2, y_2)$，$P$ 为 ${AB}$ 的中点，则有$\left\{\begin{array}{l}
             {x_1}^{2} - \dfrac{{y_1}^{2}}{2} = 1  \\
             {x_2}^{2} - \dfrac{{y_2}^{2}}{2} = 1
        \end{array}\right.$，两式作差有$(x_1+x_2)(x_1-x_2) - \dfrac{(y_1+y_2)(y_1-y_2)}{2} = 0$，根据$P$为中点有$x_1+x_2 = 2, y_1+y_2 = 2$，故$2(x_1-x_2) - (y_1-y_2) = 0$，即$k_{AB} = 2$，故$AB: y = 2x-1$

        联立直线和双曲线方程有$\left\{\begin{array}{l}
             y = 2x-1  \\
             {x}^{2} - \dfrac{{y}^{2}}{2} = 1
        \end{array}\right.$，消去$y$可得$2x^2-(2x-1)^2 = 2$，即$-2x^2+4x-3=0$，方程无解，故不存在.
    \end{jieda}
    \item 经过点$(-2,3)$,且对称轴都在坐标轴上的等轴双曲线的标准方程是\blkda[3]{$\dfrac{{y}^{2}}{5} - \dfrac{{x}^{2}}{5} = 1$}.
    \begin{jieda}
        由题意,所求双曲线为等轴双曲线,可得双曲线的方程为 ${x}^{2} - {y}^{2} = \lambda \left( {\lambda  \neq  0}\right)$ ,
        
        因为所求双曲线过点$(-2,3)$,可得 ${\left( -2\right) }^{2} - {3}^{2} = \lambda$ ,解得 $\lambda  =  - 5$，所以所求双曲线的方程为 $\dfrac{{y}^{2}}{5} - \dfrac{{x}^{2}}{5} = 1$ 
    \end{jieda}
    \item 已知双曲线的方程为 $\dfrac{{x}^{2}}{9} - \dfrac{{y}^{2}}{16} = 1$ ,点 ${F}_{1},{F}_{2}$ 是其左右焦点, ${A}$ 是圆 ${x}^{2} + {\left( y - 5\right) }^{2} = 4$ 上的一点,点 $M$ 在双曲线的右支上,则 $\left| {M{F}_{1}}\right|  + \left| {MA}\right|$ 的最小值是\blkda{$5\sqrt{2} + 4$}.
    \begin{jieda}
        由题意可得, ${c}^{2} = 9 + {16} = {25}$ ,即 $c = 5$ ,则 ${F}_{1},{F}_{2}$ 的坐标分别为 $\left( {-5,0}\right) ,\left( {5,0}\right)$，由双曲线的定义,得 $\left| {M{F}_{1}}\right|  - \left| {M{F}_{2}}\right|  = {2a} = 6$ 
        
        \begin{center}
        \includegraphics[width=0.2\textwidth]{bodyxb/src//bo_d3esa9s601uc738iqt8g_2_241_522_348_386_0.jpg}
        \end{center}

        又 ${A}$ 是圆 ${x}^{2} + {\left( y - 5\right) }^{2} = 4$ 上的点,设圆的圆心为 $C\left( {0,5}\right)$ ,半径为 $r = 2$
        
        由图可知, $\left| {CM}\right|  + \left| {M{F}_{2}}\right|  \geq  \left| {C{F}_{2}}\right|$，所以 $\left| {M{F}_{1}}\right|  + \left| {MA}\right|  = \left| {M{F}_{2}}\right|  + {2a} + \left| {MA}\right|  \geq  \left| {M{F}_{2}}\right|  + \left| {CM}\right|  + {2a} + {CA} \geq  C{F}_{2} + 6 - 2 = 5\sqrt{2} + 4$ ,
        
        当且仅当 $C$ 、 ${A}$ 、 $M$ 、 ${F}_{2}$ 四点共线( ${A}$ 、 $M$ 在 $C{F}_{2}$ 之间)时取等号,
        
        所以 $\left| {M{F}_{1}}\right|  + \left| {MA}\right|$ 的最小值为 $5\sqrt{2} + 4$
    \end{jieda}
    \item 点 $Q$ 是双曲线 $C : \dfrac{{x}^{2}}{16} - \dfrac{{y}^{2}}{4} = 1$ 上一动点,过 $Q$ 作圆 $D : {\left( x - 6\right) }^{2} + {y}^{2} = 1$ 的两条切线,切点为 $A$ ,$B$ ,则 $\left| {QA}\right|$ 的最小值为\blkda{$\dfrac{\sqrt{55}}{5}$}.
    \begin{jieda}
        由题知: 设 $Q\left( {{x}_{0},{y}_{0}}\right) ,{x}_{0} \geq  4$ ,则 $\dfrac{{x}_{0}^{2}}{16} - \dfrac{{y}_{0}^{2}}{4} = 1$，由于 $\bigtriangleup {QAD}$ 是直角三角形,且 $\left| {AD}\right|  = 1$ ,所以当 $\left| {QD}\right|$ 取得最小值时, $\left| {QA}\right|  = \sqrt{{\left| QD\right| }^{2} - 1}$ 取得最小值,则 $\left| {QD}\right|  = \sqrt{{\left( {x}_{0} - 6\right) }^{2} + {y}_{0}^{2}} = \sqrt{{\left( {x}_{0} - 6\right) }^{2} + \dfrac{{x}_{0}^{2}}{4} - 4} = \sqrt{\dfrac{5{x}_{0}^{2}}{4} - {12}{x}_{0} + {32}} = \sqrt{\dfrac{5}{4}{\left( {x}_{0} - \dfrac{24}{5}\right) }^{2} + \dfrac{16}{5}} \geq  \dfrac{4\sqrt{5}}{5}$ ,当 ${x}_{0} = \dfrac{24}{5}$ 时,等号成立,故 $\left| {QA}\right|  = \sqrt{{\left| QD\right| }^{2} - 1} \geq  \dfrac{\sqrt{55}}{5}$
        \begin{center}
            \includegraphics[width=0.5\linewidth]{bodyxb//src/shuangquxianxingzhi1B6.png}
        \end{center}
    \end{jieda}
    \item 已知点 $P$ 在焦点为 ${F}_{1}$ 、 ${F}_{2}$ 的双曲线 $\dfrac{{x}^{2}}{{a}^{2}} - \dfrac{{y}^{2}}{{b}^{2}} = 1\left( {a > 0,b > 0}\right)$ 的右支上运动, 则$\triangle P{F}_{1}{F}_{2}$ 的内切圆的圆心一定在\dotfill{(\kh{A})}
    \xx{一条直线上}{一个圆上}{一个椭圆上}{一条抛物线上}
    \begin{jieda}
        设内心为$I$，过$I$做$IH \bot x$轴，根据几何关系可知$|F_1H| - |F_2H| = |PF_1| - |PF_2| = 2a$，故$H(a, 0)$，因此$I$在直线$x = a$上.
    \end{jieda}
    \item 已知 $M\left( {-2,0}\right) \text{ 、 }N\left( {2,0}\right)$ 两点,动点 $P$ 满足条件 $\left| {PM}\right|  - \left| {PN}\right|  = 2\sqrt{2}$ ,设动点 $P$ 的轨迹为 $W$ . 
    \begin{enumerate}
        \item 求 $W$ 的方程
        \item 若 $A$ 、 $B$ 是 $W$ 上的不同两点, $O$ 是坐标原点,求 $\vv{OA} \cdot  \vv{OB}$ 的最小值
    \end{enumerate}
    \begin{jieda}
        \begin{enumerate}
            \item $\because \left| {PM}\right|  - \left| {PN}\right|  = 2\sqrt{2} < 4 = \left| {MN}\right|$
            
            由双曲线定义可知: $P$ 的轨迹是以 $M$ , $N$ 为焦点的双曲线的右支
            
            $c = 2,a = \sqrt{2},{b}^{2} = {c}^{2} - {a}^{2} = 2$
            
            因此$W$ 的方程为: $\dfrac{{x}^{2}}{2} - \dfrac{{y}^{2}}{2} = 1\left( {x > 0}\right)$
            
            \item 
            \begin{dingautolist}{172}
                \item 当直线 ${AB}$ 斜率不存在时,设直线 ${AB}$ 方程为： $x = m\left( {m > \sqrt{2}}\right)$
                此时 $A\left( {m,\sqrt{{m}^{2} - 2}}\right) ,B\left( {m, - \sqrt{{m}^{2} - 2}}\right) \therefore \vv{OA} \cdot  \vv{OB} = {m}^{2} - \left( {{m}^{2} - 2}\right)  = 2$
                \item 当直线 ${AB}$ 斜率存在时,设直线 ${AB}$ 方程为： $y = {kx} + b$
                代入双曲线方程可得： $\left( {1 - {k}^{2}}\right) {x^2} - {2kbx} - \left( {{b}^{2} + 2}\right)  = 0$
                可知上式有两个不等的正实数根 ${x}_{1},{x}_{2}\therefore \left\{  \begin{array}{l} \Delta  = 4{k}^{2}{b}^{2} + 4\left( {1 - {k}^{2}})(b^2+2) > 0\right. \\ 
                {x}_{1} + {x}_{2} = \dfrac{2kb}{1 - {k}^{2}} > 0 \\  
                {x}_{1}{x}_{2} = \dfrac{-{b}^{2} - 2}{1 - {k}^{2}} > 0 \end{array}\right.$
                解得: ${k}^{2} > 1$
                
                $\therefore \vv{OA} \cdot  \vv{OB} = {x}_{1}{x}_{2} + {y}_{1}{y}_{2} = {x}_{1}{x}_{2} + \left( {k{x}_{1} + b}\right) \left( {k{x}_{2} + b}\right)  = \left( {{k}^{2} + 1}\right) {x}_{1}{x}_{2} + {kb}\left( {{x}_{1} + {x}_{2}}\right)  + {b}^{2}$
                $= \dfrac{\left( {{b}^{2} + 2}\right) \left( {{k}^{2} + 1}\right) }{{k}^{2} - 1} + \dfrac{2{k}^{2}{b}^{2}}{1 - {k}^{2}} + {b}^{2} = \dfrac{2{k}^{2} + 2}{{k}^{2} - 1} = 2 + \dfrac{4}{{k}^{2} - 1}$
                
                由 ${k}^{2} > 1$ 得： $\therefore \vv{OA} \cdot  \vv{OB} = 2 + \dfrac{4}{{k}^{2} - 1} > 2$
            \end{dingautolist}
            综上所述, $\vv{OA} \cdot  \vv{OB}$ 的最小值为 2
        \end{enumerate}
    \end{jieda}
    \item 已知 ${F}_{1}\text{ 、 }{F}_{2}$ 为双曲线 $\dfrac{{x}^{2}}{{a}^{2}} - \dfrac{{y}^{2}}{{b}^{2}} = 1\left( {a > 0,b > 0}\right)$ 的左、右焦点,过 ${F}_{2}$ 作垂直于 $x$ 轴的直线,交此双曲线于点 $P$ ,且 $\angle P{F}_{1}{F}_{2} = {30}^{ \circ  }$ ,求此双曲线的渐近线方程.
    \begin{jieda}
        $S_{\triangle PF_1F_2} = b^2cot\ang{30} = \sqrt{3}b^2 = c \times \dfrac{b^2}{a}$，故$e = \sqrt{3}$，即$\dfrac{b}{a} = \sqrt{2}$，即渐近线方程为$y = \pm \sqrt{2}x$
    \end{jieda}
    \item 过双曲线 $y = \dfrac{k}{x}$ (常数 $k > 0$ )上任意一点 $A$ 作 ${AE}//x$ 轴,交 $y$ 轴于点 $E$ ,作 ${AF}//y$ 轴,交 $x$ 轴于点 $F$ ,得到矩形 ${AEOF}$ ,则它的面积 $S = k,k$ 是与点 $A$ 位置无关的常数. 试把这个结论推广到一般双曲线 $\dfrac{{x}^{2}}{{a}^{2}} - \dfrac{{y}^{2}}{{b}^{2}} = 1\left( {a > 0,b > 0}\right)$ ,并证明你的推广.
    \begin{jieda}
        推广结论：设 $A$ 是双曲线 $\dfrac{{x}^{2}}{{a}^{2}} - \dfrac{{y}^{2}}{{b}^{2}} = 1\left( {a > 0,b > 0}\right)$ 上任意一点,过点 $A$ 分别作渐近线 ${bx} \pm  {ay} = 0$ 的平行线 ${AE}$ 、 ${AF}$ ,并分别交渐近线于 $E$ 、 $F$ ,得到平行四边形 ${AEOF}$ ,则平行四边形 ${AEOF}$ 的面积 $S$ 是与点 $A$ 位置无关的常数.
    
        证明: 设 $A\left( {{x}_{0},{y}_{0}}\right)$ ,直线 ${AE}$ 的方程为 $y - {y}_{0} =  - \dfrac{b}{a}\left( {x - {x}_{0}}\right)$ 
        
        联立方程组 $\left\{  \begin{array}{l} y - {y}_{0} =  - \dfrac{b}{a}\left( {x - {x}_{0}}\right) \\  y = \dfrac{b}{a}x \end{array}\right.$ ,解得交点 $E\left( {\dfrac{a{y}_{0} + b{x}_{0}}{2b},\dfrac{a{y}_{0} + b{x}_{0}}{2a}}\right)$ 
        
        同理可得$F\left( {\dfrac{a{y}_{0} - b{x}_{0}}{2b}, -\dfrac{a{y}_{0} - b{x}_{0}}{2a}}\right)$

        $S = \left|\dfrac{a{y}_{0} - b{x}_{0}}{2b} \times \dfrac{a{y}_{0} + b{x}_{0}}{2a} + \dfrac{a{y}_{0} - b{x}_{0}}{2a} \times \dfrac{a{y}_{0} + b{x}_{0}}{2b}\right| = \dfrac{2|a^2y_0^2-b^2x_0^2|}{4ab} = \dfrac{2a^2b^2}{4ab} = \dfrac{ab}{2}$
        
        所以 $S = \dfrac{1}{2}a{b}$ 是与点 $A$ 位置无关的常数.
    \end{jieda}
    \item 双曲线 $\dfrac{{x}^{2}}{{a}^{2}} - \dfrac{{y}^{2}}{{b}^{2}} = 1\left( {a > b > 0}\right)$ 的实轴长 4 ,则双曲线上的一点 $\left( {4,\sqrt{3}}\right)$ 到两渐近线的距离的乘积等于\blkda{$\dfrac{4}{5}$}.
    \begin{jieda}
        设双曲线上一点为$(x_0, y_0)$，则其到双曲线的两条渐近线$y = \pm \dfrac{b}{a}x$的距离乘积为$\dfrac{\left|y_0 + \dfrac{b}{a}x_0\right|}{\sqrt{1+\dfrac{b^2}{a^2}}} \times \dfrac{\left|y_0 - \dfrac{b}{a}x_0\right|}{\sqrt{1+\dfrac{b^2}{a^2}}} = \dfrac{\left|y_0^2 - \dfrac{b^2}{a^2}x_0^2\right|}{\dfrac{c^2}{a^2}} = \dfrac{a^2b^2}{c^2}$

        实轴长 4$\Longleftrightarrow a = 2$，将$\left( {4,\sqrt{3}}\right)$代入方程，可得$b = 1$，因此$c = \sqrt{5}$，故目标式为$\dfrac{4}{5}$
    \end{jieda}
    \item 已知双曲线 $c : \dfrac{{x}^{2}}{2} - \dfrac{{y}^{2}}{{b}^{2}} = 1\left( {b > 0}\right)$ 的离心率为 $e$ ,若 $e \in  \left( {\sqrt{5},\sqrt{10}}\right)$ 则 $C$ 的焦点到一条渐近线的距离的取值范围为\blkda{$\left( {2\sqrt{2},3\sqrt{2}}\right)$}.
    \begin{jieda}
        因为 $e = \sqrt{1 + \dfrac{{b}^{2}}{{a}^{2}}} = \sqrt{1 + \dfrac{{b}^{2}}{2}} \in  \left( {\sqrt{5},\sqrt{10}}\right)$ ,所以 $b \in  \left( {2\sqrt{2},3\sqrt{2}}\right)$，而 $C$ 的焦点 $\left( {\pm c,0}\right)$ 到渐近线 ${bx} \pm  {ay} = 0$ 的距离为 $d = \dfrac{bc}{\sqrt{{a}^{2} + {b}^{2}}} = b$ .
        
        所以距离的取值范围为 $\left( {2\sqrt{2},3\sqrt{2}}\right)$ .
    \end{jieda}
    \item 已知过定点 $P\left( {0,2}\right)$ 的直线 $l$ 斜率为$k$，$l$与双曲线 $C : {x}^{2} - \dfrac{{y}^{2}}{4} = 1$ 相交于 $A,B$ 两点,问:
    \begin{enumerate}
        \item $k$ 为何值时, $A,B$ 分别在双曲线的两支上？
        \item $k$ 为何值时, $A,B$ 都在双曲线的左支上?
    \end{enumerate}
    \begin{jieda}
        双曲线 $C : {x}^{2} - \dfrac{{y}^{2}}{4} = 1$的渐近线为$y = \pm 2x$
        \begin{enumerate}
            \item 数形结合可知：$k \in (-2,2)$时满足要求.
            \item 数形结合可知：$A,B$ 都在双曲线的左支的必要条件是$k > 2$，直线$l$的方程为$y = kx + 2$，与双曲线方程联立，消去$y$可得$4x^2 - (kx+2)^2 = 4$，即$(4-k^2)x^2-4kx - 8 = 0$，$\Delta > 0 \Longleftrightarrow k^2 < 8$
            
            综上，$k \in \left( {2,2\sqrt{2}}\right)$时满足要求.
        \end{enumerate}
    \end{jieda}
    \item 双曲线 $\dfrac{{y}^{2}}{12} - \dfrac{{x}^{2}}{13} = 1$ 的一支上的三点 $A\left( {{x}_{1},{y}_{1}}\right)$ 、 ${B}\left( {\sqrt{26},6}\right)$ 、 ${C}\left( {{x}_{2},{y}_{2}}\right)$ 与焦点 $F\left( {0,5}\right)$ 的距离成等差数列. 
    \begin{enumerate}
        \item 求 ${y}_{1} + {y}_{2}$
        \item 证明:线段 ${AC}$ 的垂直平分线经过一定点
    \end{enumerate}
    \begin{jieda}
        \begin{enumerate}
            \item 根据双曲线第二定义可知$|AF| = e\left|y_1-\dfrac{a^2}{c}\right| = \dfrac{5}{2\sqrt{3}} \times \left(y_1-\dfrac{12}{5}\right)$，$|CF| = e\left|y_2-\dfrac{a^2}{c}\right| = \dfrac{5}{2\sqrt{3}} \times \left(y_2-\dfrac{12}{5}\right)$

            根据题意可知$|AF|+|CF| = 2|BF| = 6\sqrt{3}$，即$\dfrac{5}{2\sqrt{3}} \times \left(y_1+y_2-\dfrac{24}{5}\right) = 6\sqrt{3}$，故$y_1+y_2 = 12$
            
            \item $\left\{\begin{array}{l}
                 \dfrac{{y_1}^{2}}{12} - \dfrac{{x_1}^{2}}{13} = 1  \\
                 \dfrac{{y_2}^{2}}{12} - \dfrac{{x_2}^{2}}{13} = 1
            \end{array}\right.$，两式作差得到$\dfrac{(y_1+y_2)(y_1-y_2)}{12} - \dfrac{(x_1+x_2)(x_1-x_2)}{13} = 0$，即$\dfrac{y_1-y_2}{x_1-x_2} = \dfrac{12}{13} \cdot \dfrac{x_1+x_2}{y_1+y_2}$
            设$AC$中点为$M(x_0, y_0)$，则$k_{AC} = \dfrac{12}{13} \cdot \dfrac{x_0}{y_0}$

            $AC$垂直平分线方程为$y = -\dfrac{1}{k_{AC}} (x - x_0) + y_0 = -\dfrac{13y_0}{12x_0}(x-x_0)+y_0$，根据上一问可知$y_0 = 6$，故${AC}$ 的垂直平分线经过定点$\left( {0,\dfrac{25}{2}}\right)$
        \end{enumerate}
    \end{jieda}
\end{enumerate}